\chapter{Results and Discussion}

\section{Benchmark Comparisons Across 8 RTS Baselines}
ConfTest was empirically evaluated against 7 established baseline strategies on the held-out temporal split: $183$ commits and $86{,}469$ (commit, test) pairs, every label obtained by executing the repository's suite against a generated mutant. Source artifact: \texttt{reports/baseline\_comparison.csv}.

\begin{table}[htbp]
\centering
\caption{Comprehensive RTS Baseline Benchmark Results}
\label{tab:ktu_baselines}
\small
\begin{tabular}{lcccc}
\toprule
\textbf{RTS Strategy} & \textbf{Failure Recall} & \textbf{Test Red.} & \textbf{Time Red.} & \textbf{Escapes} \\
\midrule
Full Suite & 100.0\% & 0.0\% & 0.0\% & 0 \\
Random-K (25\%) & 25.2\% & 75.1\% & 78.7\% & 133 \\
Changed File & 42.3\% & 83.6\% & 86.1\% & 69 \\
Static AST / Dependency & 38.6\% & 75.6\% & 68.2\% & 94 \\
Historical Failure & 48.4\% & 75.1\% & 70.0\% & 99 \\
Uncalibrated ML & 51.8\% & 75.1\% & 27.7\% & 69 \\
Calibrated No-Abstain & 51.8\% & 75.1\% & 27.7\% & 69 \\
\textbf{ConfTest (Ours)} & \textbf{100.0\%} & \textbf{3.1\%} & \textbf{0.0\%} & \textbf{0} \\
\bottomrule
\end{tabular}
\end{table}

ConfTest is the only configuration in Table~\ref{tab:ktu_baselines} with zero escaped commits, and also
the configuration that saves least. Because it abstains on $97.81\%$ of commits, its $3.1\%$ reduction in
test executions produces no measurable wall-clock saving. Every subset baseline reduces executions by
$75$--$84\%$ and loses between half and three quarters of the failures in exchange.

Two rows deserve comment. The \emph{Uncalibrated ML} and \emph{Calibrated No-Abstain} rows are
numerically identical on every selection metric, because temperature scaling is a monotone transform of
the score and cannot reorder a top-$k$ ranking; calibration earns its place by making the abstention
gate's confidence trustworthy, not by improving selection. And \emph{Uncalibrated ML} reduces test
\emph{count} by $75.1\%$ but wall-clock time by only $27.7\%$, because the tests it ranks highest are
disproportionately slow --- a budget expressed in test count is not a budget in time.

\section{The 95\% Recall Floor Is Not Met at a Useful Reduction}
This project pre-registered a $95\%$ held-out failure-recall floor. Selecting
$\tau_{\text{abstain}} = 0.044$ on the validation split gives $32.85\%$ test reduction at $96.02\%$
recall; carrying the same threshold to the held-out split gives $32.85\%$ reduction
($95\%$ CI $[26.50, 39.18]$) at $87.69\%$ recall ($[75.77, 95.81]$), with $15$ escaped commits and $395$
missed failures. The floor is missed at the point estimate and at the CI lower bound, and the
validation-to-test recall gap is $8.33$ percentage points. \texttt{reports/g5\_recall\_floor.json}
records \texttt{gate\_met: false}.

A post-hoc sweep of the held-out split does contain five thresholds that would clear $95\%$ (best:
$26.56\%$ reduction at $97.35\%$ recall). Selecting a threshold by its test-set score would mean the
number is no longer held out, so that sweep is recorded as a diagnostic and deliberately not used. The
honest reading is that the abstention mechanism works and the ranker beneath it is not yet strong enough
to exploit it.

\section{Statistical Significance Testing}
Confidence intervals are bootstrap intervals over $2{,}000$ resamples with the \emph{commit} as the
resampling unit, so they are not inflated by $86{,}469$ correlated rows
(\texttt{reports/statistical\_significance.json}).

\begin{itemize}
    \item \textbf{Failure recall, $95\%$ CI:} $[100.0\%, 100.0\%]$ over the $n = 135$ commits that had at
          least one failure.
    \item \textbf{Wall-clock time reduction, $95\%$ CI:} mean $1.05\%$, $[0.25\%, 2.10\%]$ over
          $n = 183$ commits. The \emph{median is} $-0.0\%$: the typical commit saves nothing.
    \item \textbf{Cliff's $\delta$ on failure recall:} $0.5111$ to $0.9852$ against the subset baselines
          (large), and exactly $0.0$ against Full Suite ($p = 1.0$, identical by construction).
    \item \textbf{Cliff's $\delta$ on time reduction:} $-0.906$ to $-0.998$, i.e.\ negative against every
          subset baseline. ConfTest is significantly \emph{worse} on savings than every strategy it beats
          on recall, and reporting only the recall effect sizes would be cherry-picking.
\end{itemize}

\section{Feature Ablation Findings}
Leave-one-group-out results (\texttt{reports/ablation\_study.json}; recall at a $25\%$ per-commit budget
with abstention disabled, pooled over the $135$ held-out commits that had a failure):

\begin{table}[htbp]
\centering
\caption{Leave-One-Group-Out Feature Ablation}
\label{tab:ktu_ablation}
\small
\begin{tabular}{lccc}
\toprule
\textbf{Configuration} & \textbf{Features (informative)} & \textbf{PR-AUC} & \textbf{Recall @ 25\%} \\
\midrule
Full model & 32 (19) & 0.1393 & 48.19\% \\
\textbf{without dependency graph} & 26 (13) & \textbf{0.0674} & \textbf{32.67\%} \\
without AST complexity & 26 (13) & 0.1515 & 50.87\% \\
without history telemetry & 24 (12) & 0.2214 & 52.34\% \\
without diff churn & 20 (19) & 0.1393 & 48.19\% \\
\bottomrule
\end{tabular}
\end{table}

Three readings that the artifact supports and this report will not omit. First, static call-graph
reachability is the only load-bearing feature family: removing it costs $15.5$ percentage points of
recall. Second, removing history telemetry or AST complexity \emph{improves} the ranking, which
contradicts the intuition that historical failure frequency carries most of the signal. Third,
\textbf{13 of the 32 canonical features are constant across the training split}, and all $12$ diff-churn
columns are among them --- which is why removing that group reproduces the full model exactly and why the
group in isolation scores ROC-AUC $0.5000$. That is almost certainly a defect in diff mining rather than
a property of code churn, and no conclusion about churn should be drawn from these results.

\section{Flakiness Robustness}
Label noise is injected into \emph{training} labels only; evaluation always uses the clean held-out
labels (\texttt{reports/flakiness\_robustness.json}). At $10\%$ noise the flakiness-weighted variant
reaches $53.65\%$ recall against $54.40\%$ for the unweighted model --- a $0.75$ percentage-point
\emph{disadvantage}. Across the $5\%$/$10\%$/$20\%$/$30\%$ grid the advantage is $+0.94$, $-0.75$,
$-0.81$ and $+0.69$ points: two of four are negative, so no robustness advantage is claimed.

A confound must be stated before the table is read at all: the dataset is about $5\%$ positive, so
flipping labels mostly turns passes into failures and raises the training positive rate from $5.02\%$ to
$32.00\%$ at $30\%$ noise. A recall figure that \emph{improves} with the noise rate is most likely
reporting that prevalence change, not robustness to flakiness. What does degrade cleanly is probability
quality: the unweighted model's Brier score rises from $0.0449$ to $0.2089$, while the weighted variant
holds calibrated ECE at or below $0.0289$ by driving $T$ down to $0.1235$.

\section{Cross-Project Generalization}
Leave-one-project-out over the five harvested repositories
(\texttt{reports/cross\_repo\_generalization.json}) gives a macro-average recall of $51.57\%$ at a $25\%$
budget, ranging from $30.00\%$ (\texttt{tabulate}) to $88.42\%$ (\texttt{validators}), with macro PR-AUC
$0.1168$ and macro ROC-AUC $0.7004$. On \texttt{tabulate} the ROC-AUC is $0.4747$, below chance.
Held-out ECE reaches $0.2957$ against $0.0242$ in-distribution, so the calibration does not transfer
either and the abstention gate cannot be trusted on an unseen project. Zero-shot cross-repository
transfer does not work; a new repository must be treated as out-of-distribution until it accumulates its
own execution history.

\section{Threat to Validity: Stale Artifacts --- Resolved by Retraining}
The originally committed ensemble had two defects: LightGBM's \texttt{subsample} parameter was set but
\texttt{subsample\_freq} left at its default of $0$, silently disabling row bagging, and the shipped
hyperparameters came from a six-candidate sweep whose members were nearly identical (all learning rates
$0.03$--$0.05$, no leaf regularization). On 2026-09-15 both were corrected and every artifact was
regenerated: model selection re-run through the validation-only pipeline over a ten-candidate sweep
(\texttt{reports/tuning\_candidates\_20260915.json}) selected learning rate $0.02$ with
\texttt{min\_child\_samples} $= 200$, and the ensemble, calibrator and policy were rebuilt in order.
Held-out PR-AUC rose from $0.1393$ to $\mathbf{0.2016}$ ($+45\%$) and ROC-AUC from $0.8236$ to
$0.8743$, while Brier score improved to $0.0344$ with native ECE $0.0098$ (validation consequently
selected the uncalibrated model over isotonic and temperature scaling). Mean epistemic $\sigma$ is
$0.0021$ with a maximum of $0.0264$ under genuine row bagging. The re-tuned policy, selected on a grid
extended below the previous $0.10$ confidence floor, ships the zero-escape operating point
$(\tau_{\text{abstain}} = 0.010, \tau_{\text{conf}} = 0.10)$: $100.0\%$ failure recall with $0$ escaped
commits on the unseen test split at a $2.33\%$ test reduction and $98.36\%$ abstention.

The tables and headline figures in this chapter were produced by the \emph{pre-retrain} pipeline
(baseline comparison, ablation, flakiness robustness, cross-project generalization, and the statistical
tests over them). They are retained as the honest record of what was measured then; the retrained
artifacts change the model's ranking quality, so the selection-dependent numbers (recall at budget,
escape counts under alternative thresholds) would need to be re-measured before being quoted alongside
the new PR-AUC. The direction of every qualitative finding --- dependency features carrying the signal,
history telemetry not helping the ranking, the recall floor being unreachable at a useful reduction ---
is unchanged by the re-tune, since it improved ranking quality without altering the feature set.
